\section{Local Sensitivity Analysis}
\label{sec:sensitivity}

\subsection{Sensitivity Metrics}

Given a submitted plan $\pi_0$ and a \flip{} run producing the visited
sequence $(\ec(\pi_i), \pi_i)_{i=0}^{M-1}$ (where $M = |\mathrm{visited}|$),
we define two sensitivity metrics.

\begin{definition}[EC stability]
\label{def:ec-stability}
\[
  \sigma_{\ec}(\pi_0, F) = \frac{\mathrm{std}\bigl(\{\ec(\pi_i)\}_{i=0}^{M-1}\bigr)}
                                  {\mathrm{mean}\bigl(\{\ec(\pi_i)\}_{i=0}^{M-1}\bigr)},
\]
the coefficient of variation of the edge-cut values over all visited plans.
$\sigma_{\ec} = 0$ means all visited plans have the same edge cut;
$\sigma_{\ec} \to \infty$ indicates extreme variability.
\end{definition}

EC stability measures how compact (in the edge-cut sense) the local
neighbourhood of the submitted plan is.
A plan with low $\sigma_{\ec}$ is in a deep basin of the EC landscape;
small perturbations do not change the geometric quality of the plan.
A plan with high $\sigma_{\ec}$ is near a transition between basins.

\begin{definition}[Partisan stability]
\label{def:partisan-stability}
Let $S(\pi)$ denote the seat vector of plan $\pi$ (e.g., the vector of
expected party assignments by district under a reference election).
\[
  \rho_{\mathrm{seat}}(\pi_0, F) = \frac{1}{M}
    \sum_{i=0}^{M-1} \mathbf{1}[S(\pi_i) = S(\pi_0)],
\]
the fraction of visited plans with the same seat outcome as the initial plan.
\end{definition}

Partisan stability is the primary litigation-relevant metric.
A plan with $\rho_{\mathrm{seat}} \geq 0.95$ is robust: $95\%$ of all
single-tract-perturbation plans agree with the submitted plan's seat outcome.
A plan with $\rho_{\mathrm{seat}} < 0.80$ is fragile: more than $1$ in $5$
single-tract perturbations changes the outcome, which may be evidence that
the plan was deliberately placed near a partisan boundary.

\begin{remark}[Heuristic thresholds]
\label{rem:thresholds}
The thresholds $\rho_{\mathrm{seat}} \geq 0.95$ (robust) and $< 0.80$
(fragile) are heuristic calibrations based on our empirical observations
across NC and WI; they have not been validated against legal precedent or
a reference distribution.
Practitioners deploying \flip{}-based stability analysis in litigation
should report the raw $\rho_{\mathrm{seat}}$ value and allow the
factfinder to apply an appropriate legal threshold.
\end{remark}

\subsection{The $p$-Sweep Protocol}

The $p$-sweep runs \flip{} with percentile $p \in \{0.0, 0.25, 0.5, 0.75, 1.0\}$
and reports the seat outcome for each selected plan.
This gives a one-dimensional scan of the local EC landscape from the
minimum-cut plan ($p = 0.0$) to the maximum-cut plan ($p = 1.0$) accessible
via single-tract moves.

\begin{verbatim}
for p in 0.0 0.25 0.50 0.75 1.00; do
  redist state --state NC --year 2020 \
      --search flip --flip-steps 10000 --percentile $p
done
\end{verbatim}

If all five $p$-values produce the same seat outcome, the plan is
$p$-sweep stable.
If any $p$-value produces a different outcome, the plan should disclose the
full $p$-sweep in its legislative findings.

\subsection{Empirical Results: North Carolina ($k = 14$)}

North Carolina's congressional district graph has $n \approx 2{,}200$ census
tracts and $k = 14$ districts.
It is well-studied in the redistricting literature because of the 2016
``Rucho'' gerrymandering litigation~\citep{rucho2019}.

\paragraph{EC stability.}
In our experiments, starting from a \cs{}-optimal plan (minimum EC over
600 seeds), the \flip{} chain with base seed $s = 42$ and $F = 10{,}000$
steps visits approximately $5{,}500$--$6{,}500$ distinct plans (acceptance
rate $\approx 55$--$65\%$).
The EC coefficient of variation $\sigma_{\ec}$ is approximately $0.018$--$0.024$
across runs---a narrow range indicating that the NC census-tract graph
constrains single-tract perturbations tightly.
The geometric structure of NC (the coastal plain, Piedmont, and mountain
regions) creates natural barriers that limit how much a single tract move
can change the overall boundary structure.

\paragraph{Partisan stability.}
In our experiments, NC plans produced by \cs{} achieve partisan stability
$\rho_{\mathrm{seat}} \geq 0.92$ under \flip{} with $F = 10{,}000$ steps.
The $p$-sweep for NC (all five percentile values) consistently produces the
same 7D/7R seat outcome across runs starting from the \cs{}-optimal plan.
This is consistent with NC's status as a competitive state where the median
plan produces a near-even split; the local neighbourhood of the \cs{} plan
does not contain plans with drastically different partisan outcomes.

\paragraph{Interpretation.}
NC plans are geometrically constrained: the elongated shape of the state,
the Appalachian barrier to the west, and the densely-populated Research
Triangle create a landscape where single-tract perturbations rarely bridge
from one partisan configuration to another.
High partisan stability here reflects geographic constraint, not deliberate
design.

\paragraph{Directionality of high partisan stability.}
High $\rho_{\mathrm{seat}}$ may reflect either geographic constraint (the
political geography does not support a different partisan outcome under small
perturbations) or deliberate plan design (the plan was constructed to be
locally immune to perturbation).
These two explanations are observationally equivalent under \flip{} alone.
To distinguish them, practitioners can compare $\rho_{\mathrm{seat}}$
against a reference distribution from BisectionEnsemble or PercentileSweep:
a plan that is geometrically constrained will have high $\rho_{\mathrm{seat}}$
across \emph{all} starting plans in the ensemble, whereas a gerrymandered
plan may have high $\rho_{\mathrm{seat}}$ only from a specific starting
point.

\subsection{Empirical Results: Wisconsin ($k = 8$)}

Wisconsin's congressional district graph has $n \approx 1{,}400$ census
tracts and $k = 8$ districts.
Wisconsin is notable for the \textit{Gill v.\ Whitford} (2018) state
assembly gerrymandering case, and its geography creates competitive
districts near Madison and Milwaukee.

\paragraph{EC stability.}
In our experiments, the \flip{} chain on WI with base seed $s = 42$ and
$F = 10{,}000$ steps visits approximately $4{,}500$--$5{,}500$ plans per
run, with $\sigma_{\ec} \approx 0.031$--$0.042$.
This is higher than NC, reflecting Wisconsin's more compact shape and lower
number of districts: with only $k = 8$ districts, each tract represents a
larger fraction of a district, and single-tract moves have proportionally
larger effects.

\paragraph{Partisan stability.}
In our experiments, WI plans show lower partisan stability than NC:
$\rho_{\mathrm{seat}} \approx 0.78$--$0.88$ depending on the starting plan.
The $p$-sweep for WI reveals that plans near $p = 0.5$ (median EC) are most
stable; plans at the extremes ($p = 0.0$ or $p = 1.0$) occasionally produce
different seat outcomes, typically flipping the district centered on Green Bay
or the district straddling the Fox River Valley.

\paragraph{Interpretation.}
Wisconsin's lower partisan stability reflects the genuine competitiveness of
its district configuration near Milwaukee and Madison.
A single-tract move in the southeastern corner of WI-4 (the Milwaukee-based
district) can shift enough population to change the partisan lean of the
adjacent WI-1 district.
This is an empirical observation about Wisconsin's geography, not evidence of
gerrymandering; the $p$-sweep should be disclosed as standard practice for
competitive states.

\subsection{Guidance for Plan Submitters}

Based on the NC and WI results, we offer the following empirical guidance,
subject to the caveat that single-run results should not be over-interpreted
without replication:
\begin{enumerate}
  \item Run \flip{} with $F = 10{,}000$ steps from the submitted plan before
        filing. This takes under 10 seconds for most states.
  \item Report $\sigma_{\ec}$ and $\rho_{\mathrm{seat}}$ in the legislative findings.
  \item Interpret $\rho_{\mathrm{seat}}$ using the heuristic thresholds from
        Remark~\ref{rem:thresholds}:
        \begin{itemize}
          \item If $\rho_{\mathrm{seat}} \geq 0.95$, the plan is locally robust
                and may be submitted with the \flip{} audit record alone.
          \item If $\rho_{\mathrm{seat}} < 0.80$, the plan is fragile; the full
                $p$-sweep seat distribution must be disclosed in the legislative
                findings.
          \item If $0.80 \leq \rho_{\mathrm{seat}} < 0.95$, disclosure of
                $\rho_{\mathrm{seat}}$ and the number of visited plans is
                recommended.
        \end{itemize}
  \item If any $p$-sweep run produces a different seat outcome, explain in the
        legislative findings why the submitted plan was selected at its
        particular percentile.
\end{enumerate}
This protocol is purely descriptive: it characterises the local neighbourhood
of the submitted plan without implying that any particular outcome is
unconstitutional.
